\section{The knowledge base of the case study}\label{supp:kb}

This section gives the details of the knowledge base summarized in Section~\ref{sec:collab-kb}.

The five strands of the roundtable were identified by the author: premium principles and capital allocation;
conic finance and performance ratios; coherent risk measures and acceptance sets; axiomatic risk measures and
elicitability; and the extended-loss-ratio framework with its association. Its synthesis served as the starting
point of the panel.

The panel divided the literature into ten strands: insurance-pricing axiomatics; Choquet integrals and comonotonic
additivity; signed Choquet integrals; law-invariant coherent and convex risk measures; acceptability indices and
quasi-convex duality; return and star-shaped risk measures; elicitability; performance-ratio axiomatics; decision
theory and functional equations; and a critique of the earlier synthesis. Each strand agent returned exact
statements of if-and-only-if theorems with their hypotheses, a confidence label (high, medium or low) and a note
on its source. For example, the signed representation results of
\citet{cmmm2012,cmmm2015} were taken from their working paper versions.

There was no separate retrieval index. Instead, the strand reports were merged into a written digest of about
415{,}000 characters, which, together with the ten candidate formulations (about 245{,}000 characters), formed the
shared knowledge base of the later stages: the judge and prover agents were instructed to read these files in full
before working, and Section~\ref{sec:lit} was written from the digest. The confidence labels therefore
travelled with each statement into the stages that used it.

In addition, the author used Gemini (Google) to prepare summaries of the key works cited in
Section~\ref{sec:lit}: \citet{gerber1974}, \citet{yaari1987} and \citet{wyp1997} on premium principles;
\citet{msm1994}, \citet{cmmm2012,cmmm2015} and \citet{www2020} on signed Choquet integrals; \citet{adeh1999} and
\citet{kusuoka2001} on coherent risk measures; and \citet{chernymadan2009}, \citet{chew1983} and \citet{dekel1986} on
performance ratios and ratio representations. Each summary stated the core results of the work and their role in
this paper. The descriptions in Section~\ref{sec:lit} and the uses of these works in Sections~\ref{sec:axioms},
\ref{sec:results}, \ref{sec:necessity} and~\ref{sec:consequences} were compared with these summaries; no
discrepancy was found, and three descriptions in Section~\ref{sec:lit} were made more precise. Since the summaries
are themselves AI output, this comparison is a consistency check, not a check against the original publications.

Later, the author supplied author versions of three cited papers: \citet{www2020} (working paper of June 2019),
\citet{www2020b} (version of May 2020, with a correction of May 2023) and \citet{hlw2025} (version of June 2024).
The statements taken from them were checked against these versions: the definition of $I_h$, the
characterization theorem used in the representation step and the convexity theorem of \citet{www2020}; the
monotonicity criterion and Lipschitz bound, which are also stated there; the extension to general spaces, the
convention $\mathrm{ES}_1=\operatorname{ess\,sup}$ and the mixture representation of \citet{www2020b}; and the
axioms and the two characterization theorems of \citet{hlw2025}. No discrepancy was found. Lemma numbers that
could differ between these versions and the published articles were removed, and the description of
\citet{hlw2025} in Section~\ref{sec:lit} was rewritten from the paper.

The account of OpenAI's work and of the responses of the mathematical community in
Section~\ref{sec:std-practice} was assembled later, from two independent AI-prepared reports (see the Discussion in
Section~\ref{supp:practice}). Statements were kept where the two reports agree or where a primary source was read
directly (the GitHub repositories, the two arXiv papers and the AGMAI guidelines); OpenAI's own pages could not be
retrieved automatically; the texts of its Navier--Stokes and First Proof posts, and a localized version of its October post on academic
collaborations, were retrieved by the author in a browser, and the details of those projects are taken from them;
the October post is paraphrased rather than quoted.
